\documentclass[10pt]{article}
\usepackage[margin=1in]{geometry}
\usepackage{amsmath,microtype}
\title{Guide: A Pair of Parity Words That Forces the Trivial Cycle}
\author{Collatz proof project}
\date{September 5, 2026}
\begin{document}
\maketitle
\section*{The result}
Take any finite binary word $u$. Put an odd step at its beginning and an
even step at its end, making $1u0$. Now swap only those endpoint bits,
making $0u1$.

Lean proves that these two words cannot both describe complete returns
from natural Collatz seeds, unless $u$ is empty. In the empty case they
are simply $10$ and $01$, starting at 1 and 2 in the usual cycle.
The returning seeds do not even have to lie in the same orbit.

\section*{Why this works}
Every parity word gives an exact affine equation for its endpoint.
For a cycle, this expresses its starting seed as a correction divided by
$D=2^{\text{length}}-3^{\text{odd count}}$. On a positive natural cycle,
$D$ is positive and odd.

For our paired words, a short cancellation shows that $D$ must divide a
power of two. Therefore $D=1$. If the middle word is nonempty, reducing
$2^{\text{length}}=3^{\text{odd count}}+1$ modulo 8 gives a contradiction.
Only the empty middle survives, and its return equations give seeds 1 and 2.

\section*{How this helps the research}
The nested Knight formalization uses this cancellation to exclude a class
of evenly spaced parity patterns called high cycles. The new theorem puts
the obstruction directly into this project's actual-orbit language and
does not require a palindrome condition on the middle word.

The next task is to prove that every primitive rational mechanical pattern
contains the required pair as rotations of one period. Exact finite checks
confirm this for all reduced slopes with denominator up to 80. That finite
check does not prove the universal statement. Nonprimitive periods and
eventual tails also need explicit treatment.

\section*{The limitation}
Arbitrary parity words need not have the required pair. The word $11100$
is a concrete example with no such pair among its rotations. Its associated
cyclic value is the fraction $19/5$, not an integer counterexample.

This result is a formal integration of a known cancellation idea, not a
proof of Collatz or a claim of mathematical novelty. The full conjecture
still requires general descent, or both escape and nontrivial-cycle exclusion.

\section*{Verification}
From \texttt{lean/}, run \texttt{lake build Collatz.EndpointCycles}.
From the root, with Lean on the executable path, run
\texttt{python3 analysis/audit\_theorems.py --build}.
The three exact finite controls run with
\texttt{python3 -m unittest discover -s analysis -p test\_endpoint\_cycles.py}.
The companion technical paper gives the algebra, exact theorem scope,
provenance, and remaining formal bridge.
\end{document}
